\chapter{Objetivos}

\section{Objetivo General}

Desarrollar una aplicación gráfica interactiva en tiempo real utilizando OpenGL 3.3 y el lenguaje C++, capaz de simular una escena de lluvia mediante el uso de billboards y sistemas de partículas, incorporando además efectos ambientales, una interfaz gráfica de control y una arquitectura modular que facilite la escalabilidad y el mantenimiento del proyecto.

\section{Objetivos Específicos}

\begin{itemize}

\item Diseñar e implementar un sistema de renderizado basado en OpenGL para la representación de objetos bidimensionales y tridimensionales dentro de una escena virtual.

\item Implementar un sistema de partículas que permita generar, actualizar y reciclar miles de partículas de lluvia de forma eficiente.

\item Utilizar la técnica de billboards para representar las gotas de lluvia y los árboles, logrando una reducción significativa en el número de polígonos procesados sin afectar la percepción visual del escenario.

\item Desarrollar un terreno texturizado que sirva como superficie base para la simulación del entorno.

\item Implementar una cámara con movimiento libre sobre el plano horizontal que permita recorrer la escena de manera intuitiva.

\item Incorporar una interfaz gráfica mediante Dear ImGui que permita modificar en tiempo real parámetros como la cantidad de partículas, velocidad de la lluvia y dirección del viento.

\item Integrar un sistema de audio ambiental utilizando la biblioteca miniaudio para reproducir sonidos de lluvia y mejorar la inmersión del usuario.

\item Implementar un sistema climático simple que genere relámpagos mediante variaciones temporales de iluminación dentro de la escena.

\item Diseñar una arquitectura modular separando los diferentes componentes del sistema en módulos independientes para facilitar su mantenimiento y futura expansión.

\item Evaluar el funcionamiento del sistema verificando la correcta integración de todos los módulos y el desempeño de la aplicación bajo diferentes configuraciones de partículas.

\end{itemize}